\subsection*{给各副本分配它能独立消耗的份额}

图~\ref{fig:ch13-negative-balance}~中，A、B 都把同一 3000 当成自己的完整可用额度。改为在网络分区以前分配扣减权：A 持有 2000，B 持有 1000，合计 3000。权利表示该来源可以独立接受多少支出，每份额度只归一个来源消耗。把有限资源先分给各参与者保管的做法通常称为 escrow，在这里可理解为预留额度。

从这份共同初态出发，A 收到一项 2000 分购买，可以消耗自己的 2000。B 也收到一项 2000 分购买，却只有 1000 权利，于是等待调配或给出本次受理规则允许的拒绝。它可能还显示旧的账户总余额 3000，准入依据采用本地未消耗的专属额度~\cite{preguica2018crdt}。

这种限制会让有全局余额的购买暂时无法在 B 完成。要改变可用位置，A 可以向 B 转移权利。另取同一份初始分配：B 先消费 1000，余权为 0。A 将自己的 1000 转给 B，A 余权变为 1000。B 取得持久转移依据以后消费新取得的 1000，A 最后消费自己的 1000，总支出为 3000。

\CThirteenFigure{ch13-original13}{预分配扣减权让各来源消耗不同份额，转移先在发送方持久成立，再由接收方使用}{fig:ch13-escrow}

若 A 先通知 B 加 1000，自己却在减少以前崩溃，A 恢复后可能继续使用原 2000，B 也能使用新增 1000。这会凭空增加额度。转移先在发送来源的持久权利状态中成立，再传播接收依据，传输中断只让接收方暂时无法使用已经转出的额度。

可以用一组累计量具体表达。设 $r_i^0$ 为来源 i 的初始权利，$D_i$ 为它累计消耗的金额，$T_{ij}$ 为 i 累计转给 j 的金额。D 与 T 都从零开始，只增长。同一来源独占并串行更新自己的 $D_i$ 和整行 $T_{ij}$，更新与输入依据共同持久保存。来源 i 的可用权利为
\[
  r_i=r_i^0+\sum_jT_{ji}-\sum_jT_{ij}-D_i.
\]
转入的累计量增加可用权利，转出和消费减少它。来源只在当前可用权利至少为 x 时，接受金额 x 的消费或转出。消费把 $D_i$ 加 x，转出把对应 $T_{ij}$ 加 x。接收方读取已经成立的累计 $T_{ji}$，按同来源同目标的累计值合并，重复接收只保留较大值。接收方也可以暂时少知道某些新转入，这会使可用额度更保守。

归纳从 $\sum_i r_i^0=3000$、全部累计量为 0 开始。消费前核对 $x\le r_i$，消费以后该来源权利保持非负，总权利减少 x，总支出增加 x。转移前核对 $x\le r_i$，转移把发送方权利减 x、接收方权利加 x，总量保持。对已经成立的全体权利记录求和，转入与转出相消，因而
\[
  \sum_i r_i=3000-\sum_i D_i.
\]
每个来源的单一写入资格、自己的最新出流记录，以及更新所依赖的转入依据，都是这份论证的条件。若两个进程同时用来源 A 的 2000 分别接受两项 2000，或 A 从旧快照恢复后重用已花额度，初始分配就会被重复使用。

完整因果依据在转移链中尤其容易看清。另取三个来源，初始权利为 A=0、B=0、C=1000。C 先持久转给 A 1000，A 取得这份依据以后再转给 B 1000。B 收到的转入来自 A，但 A 能转出的原因又来自 C。传播 A 的转出时，把它所依赖的 $T_{CA}=1000$ 一并保存，合并所得状态才能解释整条权利流。

\CThirteenFigure{ch13-transfer-dependency}{A 的转出依赖先前 C 的转入，传播后继转移时共同保存这条因果依据}{fig:ch13-transfer-dependency}

对来源自己的消费和转出，恢复始终取回最新持久记录。对其他来源的记录，合并保留已经取得的完整依赖。新来源也先取得明确分配的权利，再开始消费。这样，每项允许的更新都可以回到一份实际分配与转移历史。作者综述中的有限计数器也通过分配可扣减份额保护下限，额度耗尽的副本通过拒绝或与持有剩余额度的副本同步来继续~\cite{preguica2018crdt}。

\subsection*{权利保护总量，操作身份保护同一项消费}

即使权利机制完全正确，同 S 的重复消费也可能发生。取 A 初始权利 3000，一项 S1 的价格为 1000。错误入口第一次消耗 1000，重试又消耗 1000，A 剩余 1000。全程余额非负，总支出也没有超过 3000，S1 却形成两项消费。

把 $D_A$ 从 1000 合并到 2000，只会让各副本一致记录这两项消费。要保证 S1 只消耗一次，A 的受理仲裁就共同保存 S1 的输入、原购买与消费量，后续同参恢复读取它。权利限制累计支出，完整身份则决定两次来件属于一项消费还是两项合法购买。

\CThirteenFigure{ch13-rights-identity}{同一 S1 连续消费两次也可能保持总额安全，身份绑定让重试恢复原消费}{fig:ch13-rights-identity}

外部支付再跨过一个执行边界。C1 持有固定平台范围与 \code{pay-C1}，恢复者以原输入查询或按受保护规则继续。把本地 \code{paid} 保存到集合以后，平台已经发生的 X1/X2 各占一项资金事实。集合合并会传播这些成功，平台身份及其账本证据决定实际形成了几项付款。

第十二章还为逐键保护期限结束后的恢复引入了 \code{SealPayment}。它是本书模型新增的平台能力：平台在自己的受理域关闭某身份未来的首次机会，返回原终态或可持续查询的原操作编号。已经受理的工作由平台保留到终结，恢复者据此继续核验原 C1 的成功、拒绝或未决证据。E1 关闭结束旧购买入口资格，C1/T1 则从原关联与这份平台证据继续恢复。

\subsection*{发一张券可以表达为集合，但要选对成员身份}

设活动规则是账户 1 在活动 V7 下取得一份固定权益。成员键依次包含账户 1、活动 V7 和权益类型 \code{entitlement}，授权核验后把它加入集合。两次相同授予合并为同一成员，其他账户或其他活动有自己的成员。业务只承认这一份权益，成员键据此选定区分范围。

若规则变为每完成一项购买发一张券，O1 与 O2 即使商品、金额完全相同，也应取得两张券。成员身份可以使用 \code{(campaign,\allowbreak O1)} 与 \code{(campaign,\allowbreak O2)}，每项授权生成稳定券号并保存绑定。用 \code{(account,\allowbreak campaign)} 合并会漏发第二张，用每次投递新生成的随机券号加入集合则可能多发。

\CThirteenFigure{ch13-original14}{成员键决定哪些授予可以合并，同一 O1 恢复原券，两次合法购买分别取得权益}{fig:ch13-coupons}

券已经发到外部钱包时，集合成员还要连接外部实际发券。先加本地成员、再调用钱包，调用前退出会漏掉发送。先发券、再保存成员，响应丢失可能重复发送。持久任务从原授予身份恢复，沿钱包接受的稳定身份及结果证据确定结局，再保存实际券号与原授予的关系。

同一张券的核销还有排他条件。两个离线门店都看见可用，分别接受核销，稍后把已核销按集合并成一个事实，两份商品已经交付。若产品要求一券一次核销，就把核销权预分配给单一可用位置，或在共同仲裁点决定胜者。允许离线重复兑付再事后追偿，则要据其业务规则保存每次兑付事实，核对实际交付与追偿金额。

\subsection*{退款把新的操作与原付款关联起来}

退款具有自己的业务意图。取原 X1 已成功收款 3000 分，从尚无退款的共同初态比较两条历史。历史 A 是同一项退款 RF1，申请退 1000，响应丢失后重试。历史 B 是用户分别确认 RF1、RF2，两项都退 1000。A 的预期累计退款为 1000，B 为 2000。

只保存一个出现过退 1000 的集合成员，会把 B 的两项合法退款合并。每次到达都为退款计数器新增 1000，又会把 A 的重试数成两项。退款身份与原支付 X1 的关联把这两种历史区分开，原结果让 RF1 可以恢复同一笔外部退款。

\CThirteenFigure{ch13-refund-histories}{共同从 X1 已收 3000、尚无退款开始，同一 RF1 的重试与 RF1、RF2 两项退款分别计入}{fig:ch13-refund-histories}

累计退款不超过原收款还要独立保护。取另一条初态，原收款 3000，两个来源各看到尚可退 3000，分别接受 RF1、RF2 各 2000。两项退款身份不同，总额却会超过收款。共同仲裁或预分配退款额度，把已经成功的退款和已经占用、尚未终结的退款都纳入可用金额。

来源 X1、操作 RF1 和累计金额规则分别回答退哪笔、是否重复、还能退多少。第~\ref{section:15-2}~节给出来源锁、未知占额与结果恢复协议，第~\ref{section:15-3}~节统一归纳金额关系。

回到本章起点，历史可以按业务等价关系压缩。只添加标签保存不同成员，计数器保存每源累计值，OR-set 保存活项和因果上下文，履约投影保存已核验阶段事实，扣减权保存分配、转移与消耗。购买、发券和退款要区别两项合法意图时，相关身份随之进入状态。下一章让这些事实遇到旧备份与地域切换，检查新的入口从哪里取得已经形成的绑定、资格与恢复依据。
\FloatBarrier
